\documentclass[12pt,a4paper]{article}
\usepackage{mathwizards-ingles}
\title{%
  \vspace{-1cm}
  {\Huge\bfseries Guía de Herramientas Técnicas}\\[0.3cm]
  {\Large Semana 0 — Configuración del Ecosistema de Estudio}\\[0.2cm]
  {\large\color{grissuave} Anki · FSRS · AnkiConnect · Language Reactor · Yomitan}\\[0.5cm]
  \rule{\textwidth}{1.5pt}
}
\author{}
\date{}

\begin{document}
\maketitle
\thispagestyle{empty}

\vspace{-0.5cm}

\begin{cajaalerta}[\textcolor{white}{Importante: Completar ANTES de la Sesión 1}]
Esta guía se trabaja durante la \textbf{sesión técnica preparatoria} (Semana 0), guiada paso a paso por el moderador en videollamada. Todos los estudiantes deben tener estas herramientas instaladas y configuradas \textbf{antes} de comenzar el Mes 1 del curso.
\end{cajaalerta}

% ═══════════════════════════════════════════════════════════════════════════════
\section{Anki — Sistema de Repetición Espaciada}
% ═══════════════════════════════════════════════════════════════════════════════

\begin{cajaherramienta}[¿Qué es Anki?]
Anki es un software de tarjetas de memoria (\textit{flashcards}) que utiliza algoritmos de repetición espaciada para optimizar la retención a largo plazo. Es la columna vertebral del estudio deliberado en este curso.

\medskip
\textbf{Plataformas:} Windows, macOS, Linux (gratuito) · Android (gratuito) · iOS (de pago).\\
\textbf{Descarga:} \url{https://apps.ankiweb.net/}
\end{cajaherramienta}

\subsection{Instalación}

\begin{cajapaso}[Paso 1: Descargar e Instalar]
\begin{enumerate}[leftmargin=1.2cm]
  \item Visitar \url{https://apps.ankiweb.net/}
  \item Descargar la versión para tu sistema operativo
  \item Instalar siguiendo las instrucciones del instalador
  \item Crear una cuenta gratuita en \url{https://ankiweb.net/} (para sincronización)
  \item Abrir Anki e iniciar sesión con tu cuenta
\end{enumerate}
\end{cajapaso}

\begin{cajapaso}[Paso 2: Importar el Mazo Base del Grupo]
El profesor compartirá un archivo \texttt{.apkg} con el mazo base de vocabulario de alta frecuencia.

\begin{enumerate}[leftmargin=1.2cm]
  \item Descargar el archivo \texttt{.apkg} proporcionado
  \item En Anki: \texttt{Archivo} $\rightarrow$ \texttt{Importar}
  \item Seleccionar el archivo descargado
  \item Verificar que el mazo aparezca en la pantalla principal
\end{enumerate}
\end{cajapaso}

\subsection{Configuración del Algoritmo FSRS}

\begin{cajaalerta}[Configuración Crítica — No Omitir]
El algoritmo FSRS (\textit{Free Spaced Repetition Scheduler}) reemplaza el algoritmo clásico de Anki (SM-2) y es significativamente más eficiente. Esta configuración es \textbf{obligatoria} para el curso.
\end{cajaalerta}

\begin{cajapaso}[Paso 3: Activar FSRS]
\begin{enumerate}[leftmargin=1.2cm]
  \item En Anki, ir a \texttt{Herramientas} $\rightarrow$ \texttt{Preferencias}
  \item Seleccionar la pestaña \texttt{Programación} (\textit{Scheduling})
  \item Activar la casilla \textbf{``Activar el programador FSRS''} (\textit{Enable FSRS scheduler})
  \item Hacer clic en \texttt{Aceptar}
\end{enumerate}
\end{cajapaso}

\begin{cajapaso}[Paso 4: Parámetros del Mazo]
Ir a la rueda de configuración del mazo $\rightarrow$ \texttt{Opciones}:

\medskip
\begin{tabularx}{\textwidth}{@{} l X @{}}
\toprule
\textbf{Parámetro} & \textbf{Valor Recomendado} \\
\midrule
Tarjetas nuevas por día & \textbf{5--10} (nunca más de 10) \\
Repasos máximos por día & \textbf{9999} (sin límite artificial) \\
Intervalo de graduación & \textbf{1 día, 3 días} \\
Intervalo fácil & \textbf{4 días} \\
Retención deseada & \textbf{0.90} (90\%) \\
Umbral de sanguijuela & \textbf{5 fallos} (luego suspender la tarjeta) \\
\bottomrule
\end{tabularx}
\end{cajapaso}

\begin{cajaalerta}[Regla de las ``Sanguijuelas'' (\textit{Leeches})]
Si una tarjeta falla más de 5 veces, Anki la suspenderá automáticamente. Esto es \textbf{intencional}: las sanguijuelas drenan energía cognitiva. El estudiante debe:
\begin{enumerate}[leftmargin=1.2cm, nosep]
  \item Revisar si la tarjeta está mal diseñada (¿tiene más de un elemento desconocido?)
  \item Reformularla o reemplazarla por una oración mejor (\textit{i+1})
  \item Si no logra retenerla, eliminarla sin culpa
\end{enumerate}
\end{cajaalerta}

% ═══════════════════════════════════════════════════════════════════════════════
\section{AnkiConnect — Puente de Automatización}
% ═══════════════════════════════════════════════════════════════════════════════

\begin{cajaherramienta}[¿Qué es AnkiConnect?]
AnkiConnect es un complemento (\textit{add-on}) de Anki que permite a otras aplicaciones y extensiones de navegador crear tarjetas automáticamente. Es el puente que conecta las herramientas de minería con Anki.
\end{cajaherramienta}

\begin{cajapaso}[Paso 5: Instalar AnkiConnect]
\begin{enumerate}[leftmargin=1.2cm]
  \item En Anki, ir a \texttt{Herramientas} $\rightarrow$ \texttt{Complementos} $\rightarrow$ \texttt{Obtener complementos}
  \item Ingresar el código: \textbf{2055492159}
  \item Hacer clic en \texttt{OK} y reiniciar Anki
  \item Verificar que AnkiConnect aparezca en la lista de complementos activos
\end{enumerate}
\end{cajapaso}

\begin{cajaexito}[Verificación]
Para confirmar que funciona, con Anki abierto, abrir un navegador y visitar:

\begin{cajacodigo}
http://localhost:8765
\end{cajacodigo}

Deberías ver una respuesta como: \texttt{AnkiConnect v.6}. Si aparece un error de conexión, verifica que Anki esté ejecutándose.
\end{cajaexito}

% ═══════════════════════════════════════════════════════════════════════════════
\section{Language Reactor — Inmersión con Video}
% ═══════════════════════════════════════════════════════════════════════════════

\begin{cajaherramienta}[¿Qué es Language Reactor?]
Es una extensión de navegador (Chrome/Edge) que potencia el aprendizaje de idiomas al ver contenido en YouTube y Netflix. Muestra subtítulos duales, permite pausar por oración, y facilita la exportación de vocabulario.
\end{cajaherramienta}

\begin{cajapaso}[Paso 6: Instalar Language Reactor]
\begin{enumerate}[leftmargin=1.2cm]
  \item Abrir la Chrome Web Store (o Edge Add-ons)
  \item Buscar ``Language Reactor''
  \item Instalar la extensión
  \item Abrir un video de YouTube en inglés
  \item Verificar que aparezca el panel de Language Reactor debajo del video
\end{enumerate}
\end{cajapaso}

\begin{cajapaso}[Configuración Recomendada]
\begin{tabularx}{\textwidth}{@{} l X @{}}
\toprule
\textbf{Opción} & \textbf{Configuración} \\
\midrule
Idioma de aprendizaje & English \\
Idioma nativo & \textbf{Ocultar} (no mostrar subtítulos en español) \\
Modo de reproducción & ``Auto-pause'' desactivado (para flujo natural) \\
Guardar palabras & Activar exportación a CSV (para importar a Anki) \\
\bottomrule
\end{tabularx}
\end{cajapaso}

\begin{cajaalerta}[Advertencia sobre Subtítulos]
\textbf{NO} activar subtítulos en español. El objetivo es tolerar la ambigüedad. Si necesitas apoyo, usa subtítulos en \textbf{inglés} solamente, y solo a partir del Mes 2 durante la lectura asistida.
\end{cajaalerta}

% ═══════════════════════════════════════════════════════════════════════════════
\section{Yomitan — Diccionario Emergente de Minería}
% ═══════════════════════════════════════════════════════════════════════════════

\begin{cajaherramienta}[¿Qué es Yomitan?]
Yomitan (anteriormente Yomichan) es una extensión de navegador que muestra definiciones emergentes (\textit{pop-up}) al pasar el cursor sobre palabras en inglés. Conectada con AnkiConnect, permite crear tarjetas Anki con un solo clic directamente desde cualquier página web.
\end{cajaherramienta}

\begin{cajapaso}[Paso 7: Instalar Yomitan]
\begin{enumerate}[leftmargin=1.2cm]
  \item Ir a la página oficial: \url{https://yomitan.wiki/}
  \item Instalar la extensión para Chrome/Edge/Firefox
  \item Abrir las opciones de Yomitan (clic derecho en el ícono $\rightarrow$ \textit{Options})
\end{enumerate}
\end{cajapaso}

\begin{cajapaso}[Paso 8: Cargar Diccionarios]
Yomitan necesita diccionarios descargados para funcionar:

\begin{enumerate}[leftmargin=1.2cm]
  \item Descargar diccionarios desde: \url{https://github.com/MarvNC/yomitan-dictionaries}
  \item Diccionarios recomendados para inglés:
    \begin{itemize}[nosep]
      \item \textbf{Mes 1--2:} Diccionario bilingüe Inglés$\rightarrow$Español
      \item \textbf{Mes 2--3:} \textit{Oxford Learner's Dictionary} (monolingüe simplificado)
      \item \textbf{Frecuencia:} Diccionario de frecuencia (para saber qué tan común es una palabra)
    \end{itemize}
  \item En Yomitan: \texttt{Opciones} $\rightarrow$ \texttt{Diccionarios} $\rightarrow$ \texttt{Importar}
  \item Seleccionar los archivos \texttt{.zip} descargados
\end{enumerate}
\end{cajapaso}

\begin{cajapaso}[Paso 9: Conectar Yomitan con AnkiConnect]
\begin{enumerate}[leftmargin=1.2cm]
  \item En las opciones de Yomitan, ir a la sección \texttt{Anki}
  \item Activar \textbf{``Enable Anki integration''}
  \item Configurar el tipo de nota (\textit{Note Type}) para que coincida con el formato del curso:
    \begin{itemize}[nosep]
      \item \textbf{Front:} \texttt{\{sentence\}} (la oración completa con la palabra resaltada)
      \item \textbf{Back:} \texttt{\{glossary\}} + \texttt{\{audio\}} + \texttt{\{pitch-accent\}}
    \end{itemize}
  \item Seleccionar el mazo de destino (tu mazo personal de minería)
  \item Probar: pasar el cursor sobre una palabra en una web en inglés, clic en el botón ``+'' del pop-up
  \item Verificar que la tarjeta aparezca en Anki
\end{enumerate}
\end{cajapaso}

% ═══════════════════════════════════════════════════════════════════════════════
\section{Flujo de Trabajo Integrado}
% ═══════════════════════════════════════════════════════════════════════════════

\begin{cajaexito}[Flujo Completo de Minería de Oraciones]
Una vez configuradas todas las herramientas, el flujo de minería será:

\begin{center}
\begin{tikzpicture}[
  node distance=0.6cm,
  every node/.style={
    draw, rounded corners=3pt, minimum height=1cm,
    text width=4.5cm, align=center, font=\small
  }
]
  \node[fill=herramienta!10] (a) {Leer/Ver contenido en inglés\\(YouTube, web, e-reader)};
  \node[fill=paso!10, below=of a] (b) {Encontrar una oración con\\UNA sola palabra nueva (\textit{i+1})};
  \node[fill=herramienta!10, below=of b] (c) {Yomitan: hover sobre la\\palabra $\rightarrow$ ver definición};
  \node[fill=paso!10, below=of c] (d) {Clic en ``+'' $\rightarrow$ AnkiConnect\\crea la tarjeta automáticamente};
  \node[fill=herramienta!10, below=of d] (e) {Compartir la oración en el\\canal grupal (Discord/Slack)};
  \node[fill=exito!10, below=of e] (f) {Profesor valida $\rightarrow$\\Repasar en Anki diariamente};

  \draw[->, thick] (a) -- (b);
  \draw[->, thick] (b) -- (c);
  \draw[->, thick] (c) -- (d);
  \draw[->, thick] (d) -- (e);
  \draw[->, thick] (e) -- (f);
\end{tikzpicture}
\end{center}
\end{cajaexito}

% ═══════════════════════════════════════════════════════════════════════════════
\section{Lista de Verificación Final}
% ═══════════════════════════════════════════════════════════════════════════════

\begin{cajaexito}[Checklist de la Semana 0]
Marcar cada ítem al completarlo:

\begin{enumerate}[leftmargin=1.2cm, label=$\square$]
  \item Anki instalado y cuenta creada en AnkiWeb
  \item Mazo base del grupo importado correctamente
  \item Algoritmo FSRS activado con los parámetros recomendados
  \item AnkiConnect instalado y verificado (localhost:8765)
  \item Language Reactor instalado y configurado (sin subtítulos en español)
  \item Yomitan instalado con diccionarios cargados
  \item Yomitan conectado a AnkiConnect (tarjeta de prueba creada)
  \item Discord/Slack del grupo configurado y accesible
  \item Prueba completa: minar una oración desde una web $\rightarrow$ tarjeta en Anki
\end{enumerate}
\end{cajaexito}

\vfill
\begin{center}
\rule{0.5\textwidth}{0.5pt}\\[0.3cm]
{\small\color{grissuave} Guía preparada para el curso de Inmersión en Inglés · Semana 0 · Ecosistema Refold}
\end{center}

\end{document}
